\documentclass[11pt]{article}
\usepackage[margin=0.85in]{geometry}
\usepackage[T1]{fontenc}
\usepackage{amsmath,amssymb,booktabs,array,xurl,hyperref}
\hypersetup{colorlinks=true,urlcolor=blue,linkcolor=blue}
\setlength{\emergencystretch}{3em}
\title{Noise-Shaped Target Gains Inside the Lipschitz Certificate\newline
and MnCAV Ego Geometry}
\author{Pan Dynamics Collision Avoidance Research}
\date{October 1, 2026}
\begin{document}
\maketitle

\section{Finding}
The steady-scenario accuracy losses recorded in
\path{report/SHARMA_ESTIMATOR_AUDIT_20260930.tex} came from the target-tracker
injection gains, not from the estimator structure. The gains were larger than
the certificate needs, because their shape was chosen by a criterion unrelated
to the certificate and only one scale factor was then pushed to the
certificate's feasibility edge. Selecting all three gains in one semidefinite
program, with the same Lipschitz constants and the same decay floor, lowers
the predicted radar-noise velocity error by about 32 percent and tightens the
certified ultimate bounds by 20 to 50 percent at the same time.

Rerunning the complete held-out audit with these gains and with the MnCAV ego
geometry, against a reference whose gains are again selected per scenario on
training seeds, the estimator has the lower mean error in five of six
scenarios for each of relative position, target speed and target course
(previously two, three and five of six). It remains worse in target speed at
low relative speed (0.0584 against 0.0540 m/s) and in relative position in
curved following (0.0290 against 0.0250 m); course in curved following is a
statistical tie. The price is slower convergence: the time to the convergence
tolerance rises from about 0.5 s to 0.6--1.1 s, and noise-free errors after
2 s are larger.

\section{Why the previous gains were noisy}
The previous synthesis chose a normalized shape $\ell$ by a bounded-real LMI
in the fixed region $\operatorname{Re}s\le-1$ and then searched the scalar
bandwidth $\omega_T$ for the smallest normalized ultimate bound. Three
measurements on the audited configuration explain the result.

\begin{itemize}
\item The bound objective increases with bandwidth on the feasible set, so
the search returned the smallest certifiable bandwidth: bisection of the
dissipation inequality gives feasibility edges 2.222 and 2.576 for decay
floors 0.4 and 0.8~s$^{-1}$, and the synthesized bandwidths were 2.233 and
2.591.
\item The audited errors are the radar noise passed through these gains. For
per-axis noise $\sigma=0.12/\sqrt2/\sqrt3$~m held for 0.02~s, the stationary
covariance of $\dot e=(A-Le_1^T)e-Ln$ predicts velocity errors of 0.0663,
0.0829 and 0.1406~m/s; the audited oracle target-speed RMSE values are
0.0656, 0.0849 and 0.1476~m/s.
\item With the reference's training-selected target gains $(9,27,27)$ copied
into the unchanged estimator, its mean position, speed and course errors are
lower than the scenario-tuned reference in six of six scenarios on the
held-out seeds. Those gains are not certifiable by the dissipation inequality
in any tested domain, so this run is a diagnostic and not a design.
\end{itemize}

The comparison was therefore between certified gains and gains selected
without a certificate. The remaining question was whether the certificate
itself forces the noise, or only the way the gains were chosen inside it.

\section{The change}
Let $A$ be the triple-integrator chain, $A_L=A-Le_1^T$ with physical gains
$L\in\mathbb R^3$, and $Y=PL$. The shape program minimizes $\gamma$ subject to
\begin{gather*}
\begin{bmatrix}
A_L^TP+PA_L+2\lambda_TP+\operatorname{diag}(0,\tau_qL_q^2,\tau_sL_s^2)&Pe_3&Pe_3\\
e_3^TP&-\tau_q&0\\ e_3^TP&0&-\tau_s
\end{bmatrix}\preceq0,\\
A_L^TP+PA_L+e_2e_2^T\preceq0,\qquad
\begin{bmatrix}\gamma&Y^T\\Y&P\end{bmatrix}\succeq0 .
\end{gather*}
The first block is the existing structured Lipschitz dissipation inequality
at the existing decay floor $\lambda_T=1/T_{\rm domain}$, written in physical
coordinates. The second makes $P$ an upper bound of the observability Gramian
of the velocity error, so $L^TPL\le\gamma$ bounds the stationary velocity
variance per unit white output-noise intensity. All conditions are linear in
$(P,Y,\tau_q,\tau_s,\gamma)$. The returned gains are split into a shape with
unit slowest rate and that rate; the existing bandwidth search and
free-metric certificate then run unchanged, and every downstream field keeps
its meaning.

One metric serves both inequalities, so $\gamma$ is an upper bound of the
smallest attainable variance. A local direct search over $L$ with an exact
noise objective and the same feasibility test reached a velocity error about
10 percent lower (0.0513 against 0.0572~m/s in the default domain); that
search is not adopted. The variance is a performance preference among certified
gains. It ignores the nonlinear term and the sampled realization and is not a
deterministic bound. Velocity, prediction-horizon and position-plus-velocity
weights were compared on predicted noise only; the velocity weight gave the
lowest predicted position and velocity errors in every tested domain, with
acceleration errors within 1 percent of the best, and was kept. No candidate
was selected on simulated tracking error.

The change is in \path{estimator/synthesizeTargetTrackerCertificate.m}
(\texttt{localSolveNoiseShape} replaces \texttt{localSolveNormalizedShape})
and two reported fields in \path{estimator/synthesizeNrmmObserverGains.m}.

\begin{center}\footnotesize\setlength{\tabcolsep}{4pt}
\begin{tabular}{@{}llrrrrrr@{}}
\toprule
&&&\multicolumn{2}{c}{Predicted noise RMS}&\multicolumn{3}{c}{Certified ultimate bound}\\
Declared domain&Synthesis&Gains $L$&Pos.\ (m)&Vel.\ (m/s)&Pos.\ (m)&Vel.\ (m/s)&Acc.\ (m/s$^2$)\\
\midrule
100 m, $\lambda_T=0.4$&previous&(10.22, 39.39, 53.21)&0.0188&0.0663&4.40&20.66&42.5\\
&new&(6.24, 19.31, 27.20)&0.0157&0.0452&3.54&12.46&24.6\\
80 m, $\lambda_T=0.5$&previous&(10.63, 42.64, 59.94)&0.0191&0.0704&3.48&17.24&36.8\\
&new&(6.50, 20.93, 30.23)&0.0160&0.0478&2.74&10.27&20.8\\
50 m, $\lambda_T=0.8$&previous&(11.86, 53.05, 83.18)&0.0202&0.0829&2.22&12.78&30.4\\
&new&(7.30, 26.25, 40.94)&0.0169&0.0561&1.71&7.55&16.5\\
50 m, sideslip 0.03&previous&(16.87, 107.4, 239.5)&0.0241&0.1406&3.02&23.95&81.0\\
&new&(10.21, 51.59, 116.4)&0.0200&0.0939&2.21&13.17&40.8\\
\bottomrule
\end{tabular}\end{center}
Noise values are per axis. The certified bounds remain far above the measured
errors; they are reported because the controller consumes bounds derived from
the same certificate, and they did not grow.

\section{MnCAV ego geometry}
\path{config/mncavVehicleConfig.m} copies the nominal UMN MnCAV (2021 Chrysler
Pacifica Hybrid) parameters from the Vehicle Localization repository
(\path{config/mncavVehicleParameters.json}, review of 2026-09-25, commit
\path{a1a73aad52a5d236e9936508bc763c861c38c0e8}). They describe the stock
vehicle at curb weight; the loaded mass and center of gravity of the test
vehicle are unknown. \path{config/nrmmTrackingConfig.m} now takes the ego
rear-axle distance from it (1.714395~m instead of 1.45~m). This is the only
vehicle parameter the estimator uses. The controller configuration and the
PassVeh14DOF integration runs, which replace this value by the loaded plant's,
are unchanged.

With the previous gains, the MnCAV distance changes no scenario mean of
position or speed RMSE in the fourth decimal and course by at most
0.002~deg. In the swerve scenario the ego sideslip reaches 0.053~rad, above
the default 0.05~rad branch domain, so that scenario now declares 0.06~rad.

\section{Held-out results}
The protocol is that of the audit: 54 reference gain candidates, training
seeds 101--102, test seeds 201--220, six scenarios, error after 2~s. The
campaign was rerun in full (648 training, 1,278 test, 320 oracle and 36
refinement runs in 654.6~s, plus the scenario-tuned supplement) because the
ego geometry changes the truth trajectories. The selected reference
candidates are unchanged. No run of the current estimator failed in any
variant. In the tables, P is the estimator with the previous gains (audit of
September 30), N the estimator with the new gains, and S the scenario-tuned
reference.

\begin{center}\small
\begin{tabular}{@{}lrrrrrrrrr@{}}
\toprule
&\multicolumn{3}{c}{Position RMSE (m)}&\multicolumn{3}{c}{Target speed RMSE (m/s)}&\multicolumn{3}{c}{Course RMSE (deg)}\\
Scenario&P&N&S&P&N&S&P&N&S\\
\midrule
Straight oncoming&0.0267&0.0239&0.0249&0.0666&0.0531&0.0615&0.283&0.197&0.492\\
Curved crossing&0.0267&0.0239&0.0251&0.0660&0.0490&0.0938&0.286&0.223&3.680\\
Curved following&0.0352&0.0290&0.0250&0.1476&0.0974&0.1339&0.536&0.365&0.359\\
Swerve and braking&0.0297&0.0306&0.0624&0.0763&0.0535&0.5488&0.558&0.448&6.989\\
Low relative speed&0.0291&0.0242&0.0250&0.0849&0.0584&0.0540&0.314&0.218&0.334\\
Ego weaving&0.0275&0.0235&0.0449&0.0788&0.0613&0.2757&0.357&0.307&1.438\\
\bottomrule
\end{tabular}\end{center}

Paired differences $\Delta=\mathrm{RMSE}_S-\mathrm{RMSE}_N$ over the 20 test
seeds (positive favors the estimator; 10,000 bootstrap draws of complete
seeds; pointwise intervals conditional on these scenarios):
\begin{center}\small
\begin{tabular}{@{}lrrrrrr@{}}
\toprule
&\multicolumn{2}{c}{Position (m)}&\multicolumn{2}{c}{Target speed (m/s)}&\multicolumn{2}{c}{Course (deg)}\\
Scenario&95\% interval&N lower&95\% interval&N lower&95\% interval&N lower\\
\midrule
Straight oncoming&[0.0000, 0.0021]&12&[0.0051, 0.0121]&16&[0.266, 0.321]&20\\
Curved crossing&[$-$0.0001, 0.0025]&13&[0.0401, 0.0490]&20&[3.381, 3.540]&20\\
Curved following&[$-$0.0044, $-$0.0036]&0&[0.0330, 0.0399]&20&[$-$0.018, 0.006]&9\\
Swerve and braking&[0.0304, 0.0333]&20&[0.4687, 0.5218]&20&[6.347, 6.737]&20\\
Low relative speed&[0.0005, 0.0010]&18&[$-$0.0056, $-$0.0032]&1&[0.106, 0.125]&20\\
Ego weaving&[0.0208, 0.0220]&20&[0.2102, 0.2192]&20&[1.108, 1.152]&20\\
\bottomrule
\end{tabular}\end{center}
The position advantages on the straight and curved-crossing scenarios are
about 4 percent with intervals that reach zero; they are not substantial.
Against the globally trained reference the new gains have the lower mean in
all six nominal scenarios for all three metrics. Over the 36 noisy stress
cells (six scenarios by 20 and 100~Hz sampling, doubled noise, threefold
initial offsets, radar dropout and single-track mismatch) speed and course
remain lower in every cell and position in 34; the two exceptions are the
threefold-offset cells of the two 6~s scenarios (0.0368 against 0.0357~m and
0.0362 against 0.0358~m).

With exact ego inputs and the same injection gains the two target observers
still agree (straight oncoming speed RMSE 0.053365 and 0.053356~m/s), so the
finding of the audit that the coordinate change alone does not alter accuracy
stands.

\section{What the new gains cost}
\begin{itemize}
\item \textbf{Convergence.} The slowest closed-loop rate falls from
3.0--5.0 to 1.7--2.9~s$^{-1}$, still above the decay floor. The mean time to
the convergence tolerance rises from 0.47--0.73~s to 0.63--1.13~s (the
scenario-tuned reference needs 0.87--2.16~s where it converges). Peak
transient velocity and acceleration errors fall by about 22 to 27 percent.
\item \textbf{Errors after 2~s without noise.} The residual transient is
larger: straight oncoming position 0.0008 to 0.0091~m and speed 0.0135 to
0.0294~m/s. The globally trained reference is better in four noise-free
position cells (three before), two speed cells (two before) and one course
cell (none before).
\item \textbf{Acceleration in the 6~s scenarios.} Target-acceleration RMSE
after 2~s rises from 0.132 to 0.151 and from 0.136 to 0.165~m/s$^2$ on the
straight and curved-crossing scenarios, and falls in the other four
(for example 0.448 to 0.295 in curved following).
\item \textbf{Position during the swerve.} 0.0297 to 0.0306~m.
\end{itemize}

\section{Paired benchmark against the previous synthesis}
\path{runOnlineNrmmTrackingErrorBenchmark} was run with the previous synthesis
as its paired baseline (same runtime, truth, noise draws and MnCAV geometry;
five seeds per noisy case). The new gains have lower position and velocity
RMSE in all ten noisy cases, including the two with a maneuvering target.
\begin{center}\small
\begin{tabular}{@{}lrrrrrr@{}}
\toprule
&\multicolumn{2}{c}{Position (m)}&\multicolumn{2}{c}{Velocity (m/s)}&\multicolumn{2}{c}{Acceleration (m/s$^2$)}\\
Case&previous&new&previous&new&previous&new\\
\midrule
Retained, noise&0.0283&0.0238&0.1174&0.0829&0.1788&0.1265\\
Varying target&0.0288&0.0262&0.1286&0.1105&0.3064&0.2914\\
Target lane change&0.0290&0.0274&0.1329&0.1222&0.3523&0.3487\\
Radar dropout (1 s)&0.0503&0.0439&0.1659&0.1260&0.2437&0.1812\\
Intermittent dropout&0.0299&0.0260&0.1217&0.0890&0.1845&0.1341\\
25 Hz sampling&0.0422&0.0351&0.1771&0.1234&0.2669&0.1822\\
Aggressive ego&0.0267&0.0231&0.1109&0.0921&0.1447&0.1307\\
Oncoming (8 s)&0.0255&0.0222&0.0924&0.0692&0.1256&0.1390\\
Doubled noise&0.0566&0.0475&0.2334&0.1608&0.3574&0.2417\\
Threefold offsets&0.0283&0.0250&0.1177&0.0983&0.1800&0.1719\\
Retained, no noise (1 run)&0.0001&0.0026&0.0154&0.0228&0.0065&0.0422\\
\bottomrule
\end{tabular}\end{center}
The curvature-response lag is 0.616 and 0.608~s for the varying target and
0.616 and 0.628~s for the lane change. The maximum acceleration error is
larger with the new gains in three cases (threefold offsets 0.47 to 0.80,
aggressive ego 0.36 to 0.43 and oncoming 0.30 to 0.43~m/s$^2$).

\section{Validation}
The nineteen estimator-related test classes pass 199 of 199 tests with the
change (197 of 197 before it, with the MnCAV distance). The shape test now
checks the dissipation and Gramian inequalities of the returned shape and
that the exact noise variance does not exceed the reported bound; a second
test checks that a wider target domain cannot lower that bound. The audit's
own validation passes its 103 tests and reproduces the stored
straight-oncoming seed bit for bit. The full
suite (\texttt{runtests('tests')}) passes 562 of 564 tests with no failure; two
curb-detection tests are filtered by their recorded-data assumption.

\section{Scope}
All results are synthetic, under the repository's sensor contract and exact
Sharma-model targets except where the benchmark declares model rates. The
new gains were designed after the audit's held-out results were known; the
criterion is a predicted noise variance and no simulated error was used to
choose it, but the rerun is not a preregistered confirmation. The current
estimator receives no error-based tuning, while the reference is tuned per
scenario. The certificate concerns the continuous vector field; the sampled
realization is not certified. No real-vehicle claim follows from adopting
MnCAV stock geometry.

Two losses remain. Both follow from the certificate rather than the
structure: at low relative speed the decay floor of 0.8~s$^{-1}$ keeps the
gains above the reference's uncertified choice, and in curved following the
declared 0.03~rad sideslip domain raises the velocity Lipschitz constant
from 0.80 to 2.29. Reducing the conservatism of the scalar Lipschitz
multipliers and relaxing the common-metric bound are the next steps; neither
is done here.

\section{Artifacts and reproduction}
\path{report/TARGET_NOISE_SHAPE_GAINS_20261001/} holds the complete audit
tables of the rerun (protocol, training, selections, seed-level trials,
summaries, paired intervals, oracle, numerical checks, tests, validation), the
per-domain design table, the paired benchmark tables, and the three
diagnostic tables (previous gains with MnCAV geometry, the reference's gains
in the estimator, and the floor-free certified variant).
\path{REPRODUCTION.txt} in that directory lists the commands. MATLAB R2026a
Update 3 with the existing YALMIP and SeDuMi trees; single computational
thread for the campaigns.
\end{document}
